\section{Discussion}

\subsection{What M4 Establishes}

M4 establishes a claim-boundary audit for a finite social-choice artifact. Its
technical case is deliberately narrow: the Sen24 setting under the declared
M4 encoding. Within that boundary, the paper reports a finite audit of the
repair and reporting structure, including the 48-cell paper-facing
certificate, the SAT-side \code{RepairEmpty} classification, the UNSAT-side
orbit/fiber exactness checks, and the support-truncation discipline for
shape-blind reports. These are finite-data claims under a declared artifact
interface.

The methodological contribution is the separation of layers. M4 does not make
every layer of the artifact formally verified. It makes the boundary between
layers explicit. The finite audit is one layer; the declared encoding that
defines the audited interface is another; the partial Lean-backed semantic
target for the central right condition is a third; the Python/CNF and
semantic-to-CNF correctness bridges remain future Level C obligations. The
paper's central claim is that these distinctions should be visible in the
report itself rather than hidden behind a single word such as ``verified'' or
``certified.''

\subsection{Why Claim Boundaries Matter}

The opening problem of the paper is over-reporting. A formal artifact can be
internally correct for the task it actually performed while still being used
to support a stronger claim than the evidence licenses. A finite audit can be
mistaken for a general theorem. A declared encoding can be mistaken for a
proved semantic equivalence. A partial Lean bridge can be mistaken for Python
or CNF correctness. A repair or reporting label can be treated as earned even
when no contract has justified grouping raw repairs into that report.

M4 addresses that risk by making the reportable claim explicit at each layer.
The finite certificate answers finite questions about the declared Sen24
artifact. It does not answer every semantic or implementation question around
that artifact. The Lean bridge answers a limited semantic question about the
central right-condition target. It does not certify the generated clauses. The
claim-boundary audit is therefore not a decorative qualification added after
the technical work; it is part of the technical interpretation of the work.

\subsection{The Role Of The Declared Encoding}

The declared encoding is not a weakness in the argument. It is part of the
audit discipline. A finite artifact needs an interface: a vocabulary of masks,
right-condition witnesses, obstruction shapes, cells, repair objects, and
grouped reports. By declaring that interface rather than silently treating it
as already proved equivalent to the full Lean semantics, M4 makes its claims
conditional, checkable, and bounded.

This matters especially for finite audits. Complete enumeration inside an
interface is powerful only when the interface is named. Without that
declaration, a reader could confuse exhaustive coverage of the finite table
with exhaustive coverage of all relevant semantic interpretations. M4 avoids
that move. It says what was audited, records what the audit supports, and
leaves the stronger correspondence between the artifact interface and the
Lean/Python/CNF layers as a separate future correctness obligation.

\subsection{The Role Of The Lean Bridge}

The Lean bridge gives the paper a semantic anchor without turning it into an
end-to-end verification result. The central definition,
\code{RightAtomSemantics}, packages endpoint distinctness with Lean
\code{Decisive}. The bridge also records orientation invariance for the
unordered-pair use of the artifact-level right condition. Finally,
\code{twoRightAtoms_imply_MINLIB} proves the sufficient direction from two
fixed right conditions with distinct voters to Lean \code{MINLIB}.

Those are important facts, but their scope is limited. They say that the
central right-condition target used by the paper has a Lean-level semantic
meaning and that the fixed two-right package is sufficient for Lean
\code{MINLIB}. They do not prove full selector-free equivalence to
\code{MINLIB}. They do not prove \code{\_right\_clauses} correctness, Python
encoder correctness, SAT/CNF artifact correctness, or semantic-to-CNF
correctness. The bridge therefore supports semantic discipline without
collapsing Level B semantic anchoring into Level C artifact correctness.

\subsection{Relation To M1--M3}

M4 also synthesizes lessons from the earlier layers without re-claiming them
as new results. M1 made finite evidence diagnosable by separating proof,
audit, witness, and assumption obligations. M1.5 showed that valid finite
impossibility evidence does not determine a canonical raw repair
presentation. M2 placed the general Sen theorem in the semantic bridge layer;
that theorem is background here, not a new M4 result. M3 gave conditions under
which grouped repair reports can be justified relative to a contract.

M4 combines these lessons in the declared Sen24 setting. It treats the finite
artifact as evidence, not as an all-purpose theorem; it treats repair
presentation as something that needs reporting discipline; it leaves the
general theorem in the M2 layer; and it uses the language of reportability to
ask what the artifact is entitled to say.

\subsection{Downstream Relevance}

The relevance extends beyond this small social-choice case, but only as a
methodological lesson. Formal decision-support artifacts and automated
explanation and repair systems increasingly produce claims such as
``inconsistent,'' ``impossible,'' ``explainable,'' or ``repairable.'' In
AI-mediated governance or downstream recommendation pipelines, such labels may
influence how people understand a failed rule, a rejected policy, or a
proposed relaxation. The next question is what claim the evidence actually
licenses.

M4 does not solve AI governance, prove alignment properties, or provide a
general framework for all automated systems. Its contribution is narrower: it
shows, in a finite and auditable setting, how to separate the artifact's
internal evidence from the semantic and implementation assumptions needed for
stronger public claims. That separation is relevant wherever an automated
system reports not only that a computation succeeded, but also what the result
means for explanation, repair, or downstream action.

\subsection{Closing Synthesis}

The contribution is therefore not complete verification, but disciplined
reportability. M4 shows how a finite artifact can report only within the
evidence boundary that the paper makes explicit: finite audit under a declared
encoding, a partial Lean-backed semantic target, and future correctness
obligations kept separate. That is the abstraction contract offered here. It
does not replace the missing Level C bridges; it identifies them precisely
enough that the finite report can be read without overclaiming what has been
proved.
